\chapter{Complexity — Logspace Tableau Front}
%%% generated by the blueprint pipeline from TCSlib.Complexity.CircuitComplexity.LogspaceTableauFront
\section{Overview}

Declarations of the Lean module \texttt{TCSlib.Complexity.CircuitComplexity.LogspaceTableauFront}.

\section{Declarations}

\begin{definition}[Front-tape label alphabet]
\lean{Complexity.UTab.Front.Lb}
\label{Complexity.UTab.Front.Lb}
\leanok
For natural numbers $C$ and $d$, define the finite label alphabet consisting of the
fixed labels
$rN,rN1,rN2,e0,e1,dh,dd,mh,md,xh,xd,x1,x2,th,tdl,t1,bh,bd,b1,fin,fin2,fin3$, together
with labels $ldA(i)$ for $i\in\mathrm{Fin}(C+2)$ and labels $ldD(i)$ for
$i\in\mathrm{Fin}(d+1)$.
\end{definition}

\begin{definition}[First register index]
\lean{Complexity.UTab.Front.rX}
\label{Complexity.UTab.Front.rX}
\leanok
The first register index is $0$ among five registers.
\end{definition}

\begin{definition}[Second register index]
\lean{Complexity.UTab.Front.rA}
\label{Complexity.UTab.Front.rA}
\leanok
The second register index is $1$ among five registers.
\end{definition}

\begin{definition}[Third register index]
\lean{Complexity.UTab.Front.rB}
\label{Complexity.UTab.Front.rB}
\leanok
The third register index is $2$ among five registers.
\end{definition}

\begin{definition}[Fourth register index]
\lean{Complexity.UTab.Front.rTm}
\label{Complexity.UTab.Front.rTm}
\leanok
The fourth register index is $3$ among five registers.
\end{definition}

\begin{definition}[Fifth register index]
\lean{Complexity.UTab.Front.rD}
\label{Complexity.UTab.Front.rD}
\leanok
The fifth register index is $4$ among five registers.
\end{definition}

\begin{definition}[Counter program for front computation]
\lean{Complexity.UTab.Front.prog}
\label{Complexity.UTab.Front.prog}
\leanok
\uses{Complexity.CounterProg.Instr, Complexity.UTab.Front.Lb, Complexity.UTab.Front.rA, Complexity.UTab.Front.rB, Complexity.UTab.Front.rD, Complexity.UTab.Front.rTm, Complexity.UTab.Front.rX}
For natural numbers $C$ and $d$, this defines a program with five registers and labels
$\mathrm{Lb}(C,d)$. Its instructions at each label are as follows: at $rN$ it reads the
next input bit, going to $e0$ on exhaustion or on bit $0$, and to $rN1$ on bit $1$;
$rN1$ increments register $0$ and goes to $rN2$, which prints $1$ and returns to $rN$.
At $e0$ it prints $0$ and goes to $e1$, which increments register $0$ and starts loading
register $1$. The labels $ldA(i)$ increment register $1$ while $i<C+1$, advancing $i$;
once that condition fails, the program goes to $ldD(0)$. The labels $ldD(i)$ increment
register $4$ while $i<d$, advancing $i$; once that condition fails, the program goes to
$dh$. At $dh$, it goes to $fin$ if register $4$ is zero and to $dd$ otherwise; $dd$
decrements register $4$ and goes to $mh$. At $mh$, it goes to $bh$ if register $1$ is
zero and to $md$ otherwise; $md$ decrements register $1$ and goes to $xh$. At $xh$, it
goes to $th$ if register $0$ is zero and to $xd$ otherwise; $xd$ decrements register
$0$, then $x1$ increments register $2$, and $x2$ increments register $3$ before
returning to $xh$. At $th$, it goes to $mh$ if register $3$ is zero and to $tdl$
otherwise; $tdl$ decrements register $3$, then $t1$ increments register $0$ before
returning to $th$. At $bh$, it goes to $dh$ if register $2$ is zero and to $bd$
otherwise; $bd$ decrements register $2$, and $b1$ increments register $1$ before
returning to $bh$. Finally, $fin$ prints a unary block of length equal to the value of
register $1$ and goes to $fin2$, which prints $0$ and goes to $fin3$, where the program
halts.
\end{definition}

\begin{definition}[Polynomial bound]
\lean{Complexity.UTab.Front.tb}
\label{Complexity.UTab.Front.tb}
\leanok
For fixed natural numbers $C$ and $d$, define $T(n)=(C+1)(n+1)^d$ for each natural
number $n$.
\end{definition}

\begin{definition}[Input word]
\lean{Complexity.UTab.Front.word}
\label{Complexity.UTab.Front.word}
\leanok
\uses{Complexity.UTab.Front.tb}
For natural numbers $C,d,n$, let $w$ be the Boolean word consisting of $n$ ones,
followed by a zero, followed by $T(n)$ ones and a final zero, where $T(n)=(C+1)(n+1)^d$.
\end{definition}

\begin{definition}[Five-coordinate natural number vector]
\lean{Complexity.UTab.Front.rv}
\label{Complexity.UTab.Front.rv}
\leanok
For natural numbers $a,b,c,e,f$, define a function on the five-element set
$\text{Fin}(5)$ that takes the values $a,b,c,e,f$ at indices $0,1,2,3,4$, respectively.
\end{definition}

\begin{lemma}[Value at the first index]
\lean{Complexity.UTab.Front.rv_X}
\label{Complexity.UTab.Front.rv_X}
\leanok
\uses{Complexity.UTab.Front.rX, Complexity.UTab.Front.rv}
\difficulty{1}
For $a,b,c,e,f\in\mathbb N$, the function $i\mapsto rv(a,b,c,e,f)(i)$ on $\{0,1,2,3,4\}$
takes the value $a$ at $i=0$.
\end{lemma}

\begin{lemma}[Value at the second index]
\lean{Complexity.UTab.Front.rv_A}
\label{Complexity.UTab.Front.rv_A}
\leanok
\uses{Complexity.UTab.Front.rA, Complexity.UTab.Front.rv}
\difficulty{1}
For natural numbers $a,b,c,e,f$, define $r:\text{Fin}(5)\to\bbn$ by $r(0)=a$, $r(1)=b$,
$r(2)=c$, $r(3)=e$, and $r(4)=f$. Then $r(1)=b$.
\end{lemma}

\begin{lemma}[Value at the third register]
\lean{Complexity.UTab.Front.rv_B}
\label{Complexity.UTab.Front.rv_B}
\leanok
\uses{Complexity.UTab.Front.rB, Complexity.UTab.Front.rv}
\difficulty{1}
Let $a,b,c,e,f$ be natural numbers, and let $r_B$ be the index $2$ among the five
indices $0,1,2,3,4$. The function on these indices that takes the values $a,b,c,e,f$,
respectively, has value $c$ at $r_B$.
\end{lemma}

\begin{lemma}[Value at the Fourth Position]
\lean{Complexity.UTab.Front.rv_Tm}
\label{Complexity.UTab.Front.rv_Tm}
\leanok
\uses{Complexity.UTab.Front.rTm, Complexity.UTab.Front.rv}
\difficulty{1}
Let $a,b,c,e,f$ be natural numbers, and define $r$ on $0,1,2,3,4$ by $r(0)=a$,
$r(1)=b$, $r(2)=c$, $r(3)=e$, and $r(4)=f$. Then $r(3)=e$.
\end{lemma}

\begin{lemma}[Value at the final index]
\lean{Complexity.UTab.Front.rv_D}
\label{Complexity.UTab.Front.rv_D}
\leanok
\uses{Complexity.UTab.Front.rD, Complexity.UTab.Front.rv}
\difficulty{1}
For natural numbers $a,b,c,e,f$, the function $rv(a,b,c,e,f)$ takes the value $f$ at the
index $4\in\mathrm{Fin}(5)$.
\end{lemma}

\begin{lemma}[Updating the First Register]
\lean{Complexity.UTab.Front.upd_X}
\label{Complexity.UTab.Front.upd_X}
\leanok
\uses{Complexity.UTab.Front.rX, Complexity.UTab.Front.rv}
\difficulty{4}
For natural numbers $a,b,c,e,f,v$, let $r$ be the function on five register positions
with values $a,b,c,e,f$ in order. Replacing its value at position $0$ by $v$ gives the
function with values $v,b,c,e,f$ in order.
\end{lemma}

\begin{lemma}[Updating the First Register]
\lean{Complexity.UTab.Front.upd_A}
\label{Complexity.UTab.Front.upd_A}
\leanok
\uses{Complexity.UTab.Front.rA, Complexity.UTab.Front.rv}
\difficulty{4}
For natural numbers $a,b,c,e,f,v$, let $r$ be the function on five register positions
with values $a,b,c,e,f$ in order. Replacing its value at position $1$ by $v$ gives the
function with values $a,v,c,e,f$ in order.
\end{lemma}

\begin{lemma}[Updating the Second Register]
\lean{Complexity.UTab.Front.upd_B}
\label{Complexity.UTab.Front.upd_B}
\leanok
\uses{Complexity.UTab.Front.rB, Complexity.UTab.Front.rv}
\difficulty{4}
For natural numbers $a,b,c,e,f,v$, let $r$ be the function on five register positions
with values $a,b,c,e,f$ in order. Replacing its value at position $2$ by $v$ gives the
function with values $a,b,v,e,f$ in order.
\end{lemma}

\begin{lemma}[Updating the Third Register]
\lean{Complexity.UTab.Front.upd_Tm}
\label{Complexity.UTab.Front.upd_Tm}
\leanok
\uses{Complexity.UTab.Front.rTm, Complexity.UTab.Front.rv}
\difficulty{4}
For natural numbers $a,b,c,e,f,v$, let $r$ be the function on five register positions
with values $a,b,c,e,f$ in order. Replacing its value at position $3$ by $v$ gives the
function with values $a,b,c,v,f$ in order.
\end{lemma}

\begin{lemma}[Updating the Fourth Register]
\lean{Complexity.UTab.Front.upd_D}
\label{Complexity.UTab.Front.upd_D}
\leanok
\uses{Complexity.UTab.Front.rD, Complexity.UTab.Front.rv}
\difficulty{4}
For natural numbers $a,b,c,e,f,v$, let $r$ be the function on five register positions
with values $a,b,c,e,f$ in order. Replacing its value at position $4$ by $v$ gives the
function with values $a,b,c,e,v$ in order.
\end{lemma}

\begin{lemma}[Zero-Step Run at the Same State]
\lean{Complexity.UTab.Front.goes_nil}
\label{Complexity.UTab.Front.goes_nil}
\leanok
\uses{Complexity.CounterProg.Goes, Complexity.CounterProg.run_zero, Complexity.UTab.Front.Lb, Complexity.UTab.Front.prog}
\difficulty{1}
Let $C,d,p$ be natural numbers, let $x$ be a finite bit string, let $
ho$ assign a
natural number to each of five registers, and let $\text{l}$ be any label of the front
program determined by $C$ and $d$. For every output string $o$, the program has a run of
at most zero steps from label $\text{l}$, register assignment $
ho$, and input position
$p$ to that same label, assignment, and position, appending no output.
\end{lemma}

\begin{lemma}[Reading a unary input prefix]
\lean{Complexity.UTab.Front.goes_read}
\label{Complexity.UTab.Front.goes_read}
\leanok
\uses{Complexity.CounterProg.Goes, Complexity.CounterProg.Goes.congr, Complexity.CounterProg.Goes.trans, Complexity.CounterProg.goes_inc, Complexity.CounterProg.goes_out, Complexity.CounterProg.goes_rd_true, Complexity.UTab.Front.goes_nil, Complexity.UTab.Front.prog, Complexity.UTab.Front.rX, Complexity.UTab.Front.rv, Complexity.UTab.Front.rv_X, Complexity.UTab.Front.upd_X}
\difficulty{5}
Let $C,d,n$ be natural numbers, and suppose the input is the string of $n$ ones. For
every $j\le n$, starting at the initial reading label with all five registers zero and
input position zero, the program returns to that label with the first register equal to
$j$ and all others zero, input position $j$, and exactly $j$ ones appended. The run
takes at most $3j$ steps, uniformly for every output string already printed.
\end{lemma}

\begin{lemma}[Increment-chain execution]
\lean{Complexity.UTab.Front.goes_incChain}
\label{Complexity.UTab.Front.goes_incChain}
\leanok
\uses{Complexity.CounterProg.Goes, Complexity.CounterProg.Goes.congr, Complexity.CounterProg.Goes.trans, Complexity.CounterProg.goes_inc, Complexity.UTab.Front.Lb, Complexity.UTab.Front.goes_nil, Complexity.UTab.Front.prog}
\difficulty{4}
Let $C,d,K$ be natural numbers, let $\ell_i$ be a label for each natural number $i$, and
let $r$ be one of the five registers. Suppose that for every $i<K$, the instruction at
$\ell_i$ increments register $r$ and proceeds to $\ell_{i+1}$. For any initial register
assignment and input position, for every $i\le K$, starting at $\ell_0$ reaches $\ell_i$
with register $r$ increased by $i$, all other registers unchanged, and the same input
position, appending no output, in at most $i$ steps, uniformly for every output string
already printed.
\end{lemma}

\begin{lemma}[Congruence of register assignments]
\lean{Complexity.UTab.Front.rv_congr}
\label{Complexity.UTab.Front.rv_congr}
\leanok
\uses{Complexity.UTab.Front.rv}
\difficulty{1}
If $a=a'$, $b=b'$, $c=c'$, $e=e'$, and $f=f'$ for natural numbers
$a,b,c,e,f,a',b',c',e',f'$, then the register assignments $(a,b,c,e,f)$ and
$(a',b',c',e',f')$ are equal.
\end{lemma}

\begin{lemma}[One loop iteration on the counter]
\lean{Complexity.UTab.Front.goes_xloop}
\label{Complexity.UTab.Front.goes_xloop}
\leanok
\uses{Complexity.CounterProg.Goes, Complexity.CounterProg.Goes.congr, Complexity.CounterProg.Goes.trans, Complexity.CounterProg.goes_inc, Complexity.CounterProg.goes_loop, Complexity.UTab.Front.prog, Complexity.UTab.Front.rB, Complexity.UTab.Front.rTm, Complexity.UTab.Front.rX, Complexity.UTab.Front.rv, Complexity.UTab.Front.rv_B, Complexity.UTab.Front.rv_Tm, Complexity.UTab.Front.rv_congr, Complexity.UTab.Front.upd_B, Complexity.UTab.Front.upd_Tm, Complexity.UTab.Front.upd_X}
\difficulty{5}
For natural numbers $x_0,a,b,t,e,p$, from label $\mathsf{xh}$ with register values
$(x_0,a,b,t,e)$ and input position $p$, the program reaches label $\mathsf{th}$ with
register values $(0,a,b+x_0,t+x_0,e)$ and the same input position, printing no output,
in at most $4x_0+1$ steps. This holds for every output string already printed.
\end{lemma}

\begin{lemma}[Temporary-register loop]
\lean{Complexity.UTab.Front.goes_tloop}
\label{Complexity.UTab.Front.goes_tloop}
\leanok
\uses{Complexity.CounterProg.Goes, Complexity.CounterProg.Goes.congr, Complexity.CounterProg.goes_inc, Complexity.CounterProg.goes_loop, Complexity.UTab.Front.prog, Complexity.UTab.Front.rTm, Complexity.UTab.Front.rX, Complexity.UTab.Front.rv, Complexity.UTab.Front.rv_X, Complexity.UTab.Front.rv_congr, Complexity.UTab.Front.upd_Tm, Complexity.UTab.Front.upd_X}
\difficulty{4}
Let $C,d,x_0,a,b,t,e,p$ be natural numbers, and let $x$ be any finite binary input. For
every output string $o$, there is a run of at most $3t+1$ steps of the program
determined by $C$ and $d$ that starts at the temporary-register test label with register
values $(x_0,a,b,t,e)$, input position $p$, and output $o$, and ends at the main-loop
label with register values $(x_0+t,a,b,0,e)$, the same input position $p$, and output
$o$.
\end{lemma}

\begin{lemma}[Main-loop register transfer]
\lean{Complexity.UTab.Front.goes_aloop}
\label{Complexity.UTab.Front.goes_aloop}
\leanok
\uses{Complexity.CounterProg.Goes, Complexity.CounterProg.Goes.congr, Complexity.CounterProg.Goes.trans, Complexity.CounterProg.goes_loop, Complexity.UTab.Front.goes_tloop, Complexity.UTab.Front.goes_xloop, Complexity.UTab.Front.prog, Complexity.UTab.Front.rA, Complexity.UTab.Front.rv, Complexity.UTab.Front.rv_congr, Complexity.UTab.Front.upd_A}
\difficulty{5}
For natural numbers $C,d,x_0,a,b,e,p$ and any binary input, starting at the program's
main-loop label with register values $(x_0,a,b,0,e)$ and input position $p$, the program
reaches its return-to-transfer label with register values $(x_0,0,b+a x_0,0,e)$ and the
same input position, appending no output, in at most $a(7x_0+4)+1$ steps, uniformly for
every output string already printed.
\end{lemma}

\begin{lemma}[Register transfer loop]
\lean{Complexity.UTab.Front.goes_bloop}
\label{Complexity.UTab.Front.goes_bloop}
\leanok
\uses{Complexity.CounterProg.Goes, Complexity.CounterProg.Goes.congr, Complexity.CounterProg.goes_inc, Complexity.CounterProg.goes_loop, Complexity.UTab.Front.prog, Complexity.UTab.Front.rA, Complexity.UTab.Front.rB, Complexity.UTab.Front.rv, Complexity.UTab.Front.rv_A, Complexity.UTab.Front.rv_congr, Complexity.UTab.Front.upd_A, Complexity.UTab.Front.upd_B}
\difficulty{4}
Let $C,d,x_0,a,b,e,p$ be natural numbers, and let $x$ be a finite binary input. In the
five-register counter program specified by the front-end instruction table with
parameters $C,d$, start at the branch that tests the third register, with registers
$(x_0,a,b,0,e)$, input position $p$, and any previously printed output $o$. Within at
most $3b+1$ steps, the program reaches the branch that tests the fifth register, with
registers $(x_0,a+b,0,0,e)$, the same input position $p$, and output still $o$.
\end{lemma}

\begin{lemma}[Division-loop register update]
\lean{Complexity.UTab.Front.goes_dloop}
\label{Complexity.UTab.Front.goes_dloop}
\leanok
\uses{Complexity.CounterProg.Goes, Complexity.CounterProg.Goes.congr, Complexity.CounterProg.Goes.trans, Complexity.CounterProg.goes_loop, Complexity.UTab.Front.goes_aloop, Complexity.UTab.Front.goes_bloop, Complexity.UTab.Front.prog, Complexity.UTab.Front.rD, Complexity.UTab.Front.rv, Complexity.UTab.Front.rv_congr, Complexity.UTab.Front.upd_D}
\difficulty{5}
Let $C,d,x_0,p$ be natural numbers with $1\le x_0$, and let the binary input be
arbitrary. Starting at the division-loop label with register values $(x_0,C+1,0,0,d)$
and input position $p$, the program reaches its finishing label with register values
$(x_0,(C+1)x_0^d,0,0,0)$ and the same input position, appending no output, in at most
$d\bigl(14(C+1)x_0^d+4\bigr)+1$ steps, uniformly for every output string already
printed.
\end{lemma}

\begin{definition}[Numerical bound]
\lean{Complexity.UTab.Front.bnd}
\label{Complexity.UTab.Front.bnd}
\leanok
For natural numbers $C$ and $d$, define $B=30+C+5d+14d(C+1)$.
\end{definition}

\begin{definition}[A-label with bounded index]
\lean{Complexity.UTab.Front.labA}
\label{Complexity.UTab.Front.labA}
\leanok
\uses{Complexity.UTab.Front.Lb}
For natural numbers $C,d$ and an index $i$, define an $A$-label indexed by $i$ when
$i\le C+1$, and by $0$ otherwise; the label is drawn from the family whose $A$-indices
range over $\text{Fin}(C+2)$.
\end{definition}

\begin{definition}[D-label with bounded index]
\lean{Complexity.UTab.Front.labD}
\label{Complexity.UTab.Front.labD}
\leanok
\uses{Complexity.UTab.Front.Lb}
For natural numbers $C,d$ and an index $i$, define a $D$-label indexed by $i$ when $i\le
d$, and by $0$ otherwise; the label is drawn from the family whose $D$-indices range
over $\text{Fin}(d+1)$.
\end{definition}

\begin{theorem}[Front-program computation bound]
\lean{Complexity.UTab.Front.goes_front}
\label{Complexity.UTab.Front.goes_front}
\leanok
\uses{Complexity.CounterProg.Goes, Complexity.CounterProg.Goes.congr, Complexity.CounterProg.Goes.trans, Complexity.CounterProg.goes_goto, Complexity.CounterProg.goes_halt, Complexity.CounterProg.goes_inc, Complexity.CounterProg.goes_out, Complexity.CounterProg.goes_pr, Complexity.CounterProg.goes_rd_end, Complexity.UTab.Front.bnd, Complexity.UTab.Front.goes_dloop, Complexity.UTab.Front.goes_incChain, Complexity.UTab.Front.goes_read, Complexity.UTab.Front.labA, Complexity.UTab.Front.labD, Complexity.UTab.Front.prog, Complexity.UTab.Front.rA, Complexity.UTab.Front.rD, Complexity.UTab.Front.rX, Complexity.UTab.Front.rv, Complexity.UTab.Front.rv_A, Complexity.UTab.Front.rv_D, Complexity.UTab.Front.rv_X, Complexity.UTab.Front.tb, Complexity.UTab.Front.upd_A, Complexity.UTab.Front.upd_D, Complexity.UTab.Front.upd_X, Complexity.UTab.Front.word}
\difficulty{5}
For natural numbers $C,d,n$, on input consisting of $n$ ones and starting at the initial
reading label with all five registers zero, there exist final register values and a
final input position such that the program halts and appends the word consisting of $n$
ones, then a zero, then $(C+1)(n+1)^d$ ones, then a zero. This run relation holds
uniformly for every output string already printed, with a step bound of
$\bigl(30+C+5d+14d(C+1)\bigr)(n+1)^{d+1}$.
\end{theorem}

\begin{theorem}[Bounded computation of the output word]
\lean{Complexity.UTab.Front.run_word}
\label{Complexity.UTab.Front.run_word}
\leanok
\uses{Complexity.CounterProg.St, Complexity.CounterProg.init, Complexity.CounterProg.run, Complexity.UTab.Front.Lb, Complexity.UTab.Front.bnd, Complexity.UTab.Front.goes_front, Complexity.UTab.Front.prog, Complexity.UTab.Front.rv, Complexity.UTab.Front.word}
\difficulty{4}
For every $n\in\mathbb N$, on input consisting of $n$ consecutive $1$ bits, the program
starting at its initial label halts after some $t\leq (30+C+5d+14d(C+1))(n+1)^{d+1}$
steps and has output $1^n0 1^{(C+1)(n+1)^d}0$.
\end{theorem}
